\documentclass[11pt]{article}

\usepackage[margin=25mm]{geometry}

\usepackage{amsmath,amssymb,fontspec,unicode-math}

\IfFontExistsTF{DejaVu Serif}{\setmainfont{DejaVu Serif}}{\setmainfont{Latin Modern Roman}}

\IfFontExistsTF{Latin Modern Math}{\setmathfont{Latin Modern Math}}{\setmathfont{DejaVu Math TeX Gyre}}

\usepackage{xurl,hyperref}

\setlength{\parindent}{0pt}\setlength{\parskip}{6pt}

\begin{document}

\title{The Future Unmarked: Watermark Removal in AI-Generated Images via Next-Frame Prediction}

\author{ELVIA preliminary AI analysis}\date{21 October 2026}\maketitle

Preliminary assessment: Minor.\par

Independent scientific review remains pending. Author adoption has not been requested.

Independent horizontal and vertical searches need not maximize the combined latent-distance surrogate. A separate exact score example shows why a low detection rate at one operating point does not establish statistical indistinguishability. Neither finding remeasures the reported image attack.

\section*{Review scope}

Claim-focused reading of the original paper and relevant author-released passages; direct mathematical or static-source reasoning. No author program, earlier verifier, model or benchmark was executed during this check.

\section*{Problem 1: Independent axis maxima can cancel in the combined warp}

Source: Section 3.2, Eqs.(1)–(2), PDF p.5; utils.py, max\_warp\_latents

Assessment: supported; impact: minor

\subsection*{Source claim (paraphrase)}

The displayed target maximizes joint latent distance. The release maximizes \(L(d_x,0)\) and \(L(0,d_y)\) separately and combines their offsets.

\subsection*{Analysis}

For the accepted four-channel \(64\times64\) latent \(z_{c,i,j}=(-1)^{i+j}\), the released \(512\)-pixel flow grid induces latent source coordinate \(j+d_x/8-1/16\), with nearest-neighbor reflection. Offset \(-27\) flips parity at \(61\) of \(64\) positions per axis. Hence \(L(-27,0)=L(0,-27)=1.90625\), but the selected combination gives \[L(-27,-27)=4(61/64)(3/64)=0.1787109375.\] Offset \(20\) flips only \(2\) positions, so \[L(-27,20)=2\{(61\cdot62+3\cdot2)/4096\}=1.849609375.\] Both pairs lie within the source’s own candidate product and \([-40,40]^2\).

\subsection*{Finding and scope}

The accepted-input heuristic need not solve the joint surrogate objective. Reachability of this checkerboard by natural-image inversion, and any resulting attack-quality effect, remain unestablished.

\subsection*{Proposed correction}

Joint finite enumeration returns a maximizer over its declared candidate product. It does not solve unrestricted continuous flow optimization or prove better watermark removal.

A proposed correction is retained where stated; independent verification and author adoption remain unestablished.

\section*{Problem 2: One low-FPR operating point does not imply indistinguishability}

Source: Section 4.1, PDF p.7; Section 4.2 RQ1, Table 1 and Figure 3, PDF p.8

Assessment: supported; impact: minor

\subsection*{Source claim (paraphrase)}

The discussion infers indistinguishability from shifted detector distances and low \(\mathrm{TPR}\) at \(1\%\) false-positive rate.

\subsection*{Analysis}

Let lower distance indicate a watermark. Take negative scores \(0,1,\ldots,99\) and positive scores consisting of seven copies of \(-0.5\) and ninety-three copies of \(7.5\). Threshold \(0.5\) gives \(\mathrm{FPR}=0.01\) and \(\mathrm{TPR}=0.07\). Yet the pairwise AUC is \[(7\cdot100+93\cdot92)/10000=0.9256.\] These are explicitly constructed scores, not the paper’s unpublished image scores. Repository issue 2, opened by XiaoxiKong on 27 December 2025, independently raised the high-AUC concern; its empirical measurements were not reproduced here.

\subsection*{Finding and scope}

The operating-point result remains meaningful, but does not establish equality of score distributions or absence of all watermark information.

\subsection*{Proposed correction}

State that the evaluated detector’s sensitivity falls at the threshold calibrated to \(1\%\) false-positive rate; retain the distinction from distributional indistinguishability.

A proposed correction is retained where stated; independent verification and author adoption remain unestablished.

\section*{Problem 3: Distributional image-quality scores do not certify each image}

Source: Image-quality evaluation, PDF pp.8–9 and Appendix quality tables

Assessment: not\_issue; impact: minor

\subsection*{Source claim (paraphrase)}

FID and aggregate CLIP scores assess distributional quality and prompt alignment.

\subsection*{Analysis}

Permuting generated images within a prompt-compatible collection preserves distributional FID and aggregate semantic scores while changing which source image receives which output. Thus those summaries alone cannot certify instance-wise preservation. The paper’s reported scores themselves are not contradicted.

\subsection*{Finding and scope}

This is a metric-scope clarification; no separate image-level defect is established.

\section*{Limits}

Preliminary AI assessment; independent human scientific adjudication remains pending.\par

Analytic witnesses and static release checks do not remeasure the published benchmark results.\par

\section*{Source identity}

NeurIPS 2025 official proceedings PDF (retained local copy)

\url{https://proceedings.neurips.cc/paper_files/paper/2025/file/b48405307576689bea14d221ccba73ae-Paper-Conference.pdf}

Reviewed PDF SHA-256: \nolinkurl{fb07425e8f44b564920b8455ec48e8819d28a8671d0dc03f09167eda161e2eea}

The SHA-256 identifies the exact retained PDF inspected in this rapid source/argument check. Original supplied report bytes are retained locally. No unambiguous original PDF retrieval timestamp was recovered; file modification time was not substituted.

\end{document}
